\chapter{CAM and BIM: Machining and Architecture}

These two workbenches are aimed at professional use: CAM generates numerical-control toolpaths, and BIM builds
building information models. Their commands all come from FreeCAD's official Python modules and behave as
they do on the desktop, but the mobile port makes its own trade-offs in each.

\section{CAM: from model to G-code}

\subsection{A typical workflow}

\begin{enumerate}
  \item Open \menupath{More ▾ → CAM};
  \item \ui{Job} (\cmd{CAM\_Job}): creates a machining job from the selected solids (if nothing is selected,
        all visible solids are used); the default template sets up the job coordinate system and output settings;
  \item The job panel has three pages: \ui{Setup}, \ui{Output}, and \ui{Tools};
  \item Tools: \ui{Tool controller} → \ui{New tool} / \ui{Load from library} (the tool library ships with the app);
  \item Operations: \ui{Profile}, \ui{Pocket shape}, \ui{Mill facing},
        \ui{Helix}, \ui{Adaptive}, \ui{Slot}, \ui{Drilling},
        \ui{Engrave}, \ui{Pocket 3D};
  \item Dress-up: Dogbone, RampEntry, Tags, LeadInOut,
        array/mirror, boundary, Z correction, axis mapping, drag knife, plunge milling, and more;
  \item \ui{Sanity}: checks whether the toolpath is empty and whether feed rates or spindle speeds are missing;
  \item \ui{Post-process} (\cmd{CAM\_Post}): choose a post-processor (fanuc, grbl, linuxcnc, centroid, estlcam, and others),
        generate the G-code, and save it through the system file picker.
\end{enumerate}

\subsection{Scope and limitations}

\begin{itemize}
  \item The post-processor scripts ship with the app and cover the common 3-axis and 4-axis controllers;
  \item \textbf{There is no material-removal simulation}: the CAMSimulator that
        \cmd{CAM\_Simulator}/\cmd{CAMotics} depend on is not packaged, so you can only inspect the toolpath
        curves in the 3D view and run the basic checks with \ui{Sanity};
  \item Complex 3D toolpaths (Adaptive, Pocket 3D) are slow to compute on a phone; preview them first with a coarse tolerance.
\end{itemize}

\begin{practice}[Exercise: Mill a pocket]
Create a block with a filleted pocket → CAM job → New tool: φ6 flat-end mill → choose the pocket floor for the Pocket operation →
set the step-over and step-down → Sanity → post-process (grbl) → save \file{pocket.nc}.
\end{practice}

\section{BIM: buildings and IFC}

The BIM workbench mirrors FreeCAD's BIM module (131 commands), including:

\begin{itemize}
  \item \textbf{Architectural elements}: Site, Building, Level, Space,
        Wall, CurtainWall, Column, Beam, Slab, Door,
        Window, Covering, Pipe, Stairs, Roof,
        Panel, Frame, Fence, Truss, Equipment, Rebar;
  \item \textbf{Structural elements}: columns, beams, and slabs; structural systems; rebar; trusses; frames; curtain walls;
  \item \textbf{Materials}: Material, MultiMaterial, BIM Material;
  \item \textbf{Annotations and sheets}: dimensions, text, axis grids, sections, 2D views, sheets;
  \item \textbf{Management}: project setup, project manager, IFC manager, layers, materials, preflight checks.
\end{itemize}

\subsection{IFC import and export}

\begin{itemize}
  \item IfcOpenShell 0.8.4 ships with the app;
  \item \textbf{Open IFC} (\file{.ifc}): opens the file as a project document in NativeIFC mode;
  \item \textbf{Import IFC}: inserts the model into the current document as objects (on mobile the default is the
        ``solid shape'' strategy and all children are imported, because a phone cannot expand and load them
        on demand the way the desktop can);
  \item \textbf{Export IFC}: to \file{.ifc} or \file{.ifcJSON}; there are also commands such as ``Create IFC project'',
        ``Save IFC'', ``Expand'', ``Update'', and ``Compare'';
  \item Note: IFC uses the full schema, so both files and memory use are large; on a phone, import partial or
        simplified models first.
\end{itemize}

\section{Robot}

The Robot workbench is for kinematics demonstrations and trajectory planning:

\begin{itemize}
  \item The default robot is a KUKA IR500 (the DH parameters are built in; the VRML model is not packaged,
        so the robot is drawn from its kinematic chain);
  \item You can create trajectories and waypoints (TCP or picked points), generate a trajectory from an edge,
        dress up or compound trajectories, home the robot, and set the tool shape;
  \item \ui{Simulation} runs 60 frames using inverse kinematics;
  \item You can export a KUKA \file{.src} program (through the system save dialog).
\end{itemize}

\section{Other specialist workbenches}

\ui{Inspection} and \ui{Material} were introduced in Chapter 8, and \ui{Reverse Engineering} is there too.
Like CAM and BIM, these workbenches cover the full command set, but the mobile screen and memory limit them
to viewing, checking, and small edits; for large projects it is still better to finish on the desktop.
